\section{Kongruencje.}
Przejdziemy teraz do nieco bardziej abstrakcyjnego pojęcia,
	tak zwanych kongruencji.
Mówiąc nieformalnie,
	są to relacje równoważności na modelach algebr,
	które są spójne ze strukturą tych modeli.

\begin{definition}
	Niech $\mathfrak{A} = (\Sigma, E)$ będzie algebrą,
		$A$ modelem $\mathfrak{A}$,
		a $\sim$ dowolną relacją równoważności na $A$.
	Wtedy $\sim$ jest \textbf{kongruencją},
		jeśli dla każdego symbolu $f^n \in \Sigma$,
			reprezentującej go funkcji $f: A^n \to A$ i
	\begin{equation}
		x_1, \dots, x_n,\; y_1, \dots, y_n \in A
	\end{equation}
	 zachodzi
	\begin{equation}
		x_1 \sim y_1, \dots, x_n \sim y_n \implies f(x_1, \dots, x_n) \sim f(y_1, \dots, y_n).
	\end{equation}
\end{definition}

Czyli kongruencja jest relacją równoważności,
	która zachowuje równoważność przy nakładaniu funkcji reprezentujących symbole algebry.

\begin{example}
	Rozważmy monoid ciągów znaków z konkatenacją.
	Relacja, która identyfikuje dwa ciągi, jeśli ich długość jest taka sama,
		jest kongruencją.
\end{example}

\begin{example}
	Rozważmy $\mathbb{Z}$ jako pierścień z dodawaniem i mnożeniem.
	Dla dowolnej liczby $n \in \mathbb{N}$ relacja $\sim$ zdefiniowana przez warunek
	\begin{equation}
		\forall_{x, y \in \mathbb{Z}} \; x \sim y \iff x \mod n = y \mod n
	\end{equation}
	jest kongruencją.
\end{example}

\begin{theorem}\label{th:cong-model}
	Niech $\mathfrak{A} = (\Sigma, E)$ będzie algebrą,
		$A$ jej modelem i $\sim$ kongruencją na $A$.
	Wtedy $A/_\sim$ też jest modelem $\mathfrak{A}$.
\end{theorem}
\begin{proof}
	Musimy pokazać,
		że istnieją interpetacje symboli z $\Sigma$ jako funkcje nad $A$ z odpowiednimi
		arnościami,
		i że spełnione są przy tym wszystkie równania z $E$.
	Niech $f^n \in \Sigma$ będzie dowolnym symbolem arności $n$,
		któremu odpowiada funkcja $f$ w $A$.
	Zdefiniujemy teraz odpowiadającą mu funkcję $f_\sim$ na $A/_\sim$ jako
	\begin{equation}
		f_\sim([x_1]_\sim, \dots, [x_n]_\sim) = [f(x_1, \dots, x_n)]_\sim.
	\end{equation}
	Musimy jednak pokazać,
		że definicja $f_\sim$ jest niezależna od wybranych reprezentantów klas $[x_1]_\sim, \dots, [x_n]_\sim$.
	Weźmy więc dowolne
	\begin{equation}
		x_1, \dots, x_n,\;y_1, \dots, y_n \in A
	\end{equation}
		takie,
		że $x_i \sim y_i$ dla każdego $i = 1, \dots, n$,
		czyli $[x_i]_\sim = [y_i]_\sim$.
	Mamy teraz
	\begin{equation}
		[f(x_1, \dots, x_n)]_\sim = [f(y_1, \dots, y_n)]_\sim,
	\end{equation}
	ponieważ z definicji kongruencji mamy
	\begin{equation}
		x_1 \sim y_1, \dots, x_n \sim y_n \implies f(x_1, \dots, x_n) \sim f(y_1, \dots, y_n).
	\end{equation}
	Oznacza to,
		że definicja $f_\sim$ nie jest zależna od wybranych reprezentantów klas $[x_1]_\sim, \dots, [x_n]_\sim$,
	czyli poprawnie zdefiniowaliśmy funkcje na $A/_\sim$ odpowiadające symbolom z $\Sigma$.

	Dowód,
		że funkcje te spełniają równania z $E$,
		pozostawiamy jako zadanie dla czytelnika (zad.~\ref{ex:first-cong}--\ref{ex:last-cong}).
\end{proof}

\begin{definition}
	\textbf{Jądrem} homomorfizmu $h: A \to B$ nazywamy relację określoną wzorem
	\begin{equation}
		x \equiv_h y \iff h(x) = h(y).
	\end{equation}
\end{definition}

Definicja ta nie do końca zgadza się z tym,
	czego mogliśmy się kiedyś dowiedzieć.
Czy jądrem nie powinien być zbiór miejsc zerowych?
Dlaczego tutaj jest to relacja?
Na te pytania odpowiemy niedługo.

\begin{lemma}
	Niech $\mathfrak{A} = (\Sigma, E)$ będzie algebrą, a $A, B$ jej modelami.
	Niech $h: A \to B$ będzie homomorfizmem.
	Wtedy jego jądro $\equiv_h$ jest kongruencją.
\end{lemma}
\begin{proof}
	Fakt,
		że $\equiv_h$ jest relacją równoważności,
		pozostawiamy dla czytelnika (zad.~\ref{ex:ker-equiv}).
	Zakładając jego prawdziwość pokażemy,
		że $\equiv_h$ jest kongruencją.

	Niech $f^n \in \Sigma$ będzie dowolnym symbolem arności $n$,
		a $f_A, f_B$ odpowiadającymi mu funkcjami na odpowiednio $A, B$.
	Weźmy dowolne
	\begin{equation}
		x_1 \dots, x_n,\;y_1, \dots, y_n \in A
	\end{equation}
	spełniające
	\begin{equation}
		x_1 \equiv_h y_1, \dots, x_n \equiv_h y_n
	\end{equation}
	Chcemy pokazać, że
	\begin{equation}
		f_A(x_1, \dots, x_n) \equiv_h f_A(y_1, \dots, y_n).
	\end{equation}
	Z założenia i definicji jąda mamy
	\begin{equation}
		h(x_1) = h(y_1), \dots, h(x_n) = h(y_n),
	\end{equation}
	czyli
	\begin{equation}
		f_B(h(x_1), \dots, h(x_n)) = f_B(h(y_1), \dots, h(y_n)),
	\end{equation}
	co z faktu, że $h$ jest homorfizmem daje
	\begin{equation}
		h(f_A(x_1, \dots, x_n)) = h(f_A(y_1, \dots, y_n)),
	\end{equation}
	co jest z definicji równoważne udowadnianej przez nas własności.
\end{proof}

Skoro jądro homomorfizmu $h: A \to B$ jest relacją równoważności,
	to dzieli ono $A$ na klasy równoważności.
Jeśli algebra,
	której $A$ i $B$ są modelami ma jakiś element $e$,
	który nazwać możemy jednością
	(czyli np. operujemy na monoidach, grupach, przestrzeniach liniowych),
	to jądrem możemy też nazywać klasę abstrakcji zawierającą ten element
	(co oznaczać będziemy przez $\mathrm{Ker}\,h$).
Takie spojrzenie jest zgodne z tym,
	że jądro jest zbiorem ,,miejsc zerowych'' homomorfizmu:
\begin{align}
	[e]_{\equiv_h} &= \\
		&= \{x \in A | x \equiv_h e \} \\
		&= \{x \in A | h(x) = h(e) \} \\
		&= \{x \in A | h(x) = e \}\\
		&= \mathrm{Ker}\,h.
\end{align}

Pojęciem dualnym do jądra jest obraz homomorfizmu $\mathrm{Im}\,h$.
Jest to nic innego,
	jak zbiór wartości.
Obraz oczywiście nie musi być równy przeciwdziedzinie.
Jeśli jest,
	to $h$ jest epimorfizmem.
Jak się okazuje,
	obrazy,
	podobnie jak jądra (zad.~\ref{ex:ker-submodel}) są podmodelami.

\begin{lemma}
	Niech $h: A \to B$ będzie homomorfizmem.
	Zbiór wartości $\mathrm{Im}\,h$ jest wtedy podmodelem $B$.
\end{lemma}
\begin{proof}
	Musimy pokazać,
		że funkcje na $B$ odpowiadające symbolom z $\Sigma$ możemy ograniczyć do $\mathrm{Im}\,h$,
		czyli że dla argumentów z $\mathrm{Im}\,h$ przyjmują wartości z $\mathrm{Im}\,h$.
	Niech więc $f^n \in \Sigma$ będzie $n$-arnym symbolem,
		a $f_A, f_B$ odpowiadającymi mu funkcjami na odpowiednio $A, B$.
	Niech $y_1, \dots, y_n \in \mathrm{Im}\,h$.
	Istnieją więc $x_1, \dots, x_n \in A$,
		że
	\begin{equation}
		y_1 = h(x_1), \dots, y_n = h(x_n).
	\end{equation}
	Mamy więc
	\begin{equation}
		f_B(y_1, \dots, y_n) = f_B(h(x_1), \dots, h(x_n)) = h(f_A(x_1, \dots, x_n)),
	\end{equation}
	co znaczy,
		że $f_B(y_1, \dots, y_n) \in \mathrm{Im}\,h$.
\end{proof}

Gotowi jesteśmy teraz,
	by zapostulować i udowodnić centralne twierdzenie tego skryptu,
	które mówi,
	jak dokładnie działają homomorfizmy.

\begin{theorem}\label{th:isomorphism}
	Niech $h: A \to B$ będzie homomorfizmem.
	Wtedy
	\begin{equation}
		h_*: A/_{\equiv_h} \to \mathrm{Im}\,h,
	\end{equation}
	zdefiniowany jako
	\begin{equation}
		h_*([x]_{\equiv_h}) = h(x),
	\end{equation}
	jest izomorfizmem.
\end{theorem}
\begin{proof}
	Definicja $h_*$ jest poprawna,
		t.j. nie zależy od wyboru reprezentanta klasy.
	Wynika to bezpośrednio z definicji jądra.
	Jeśli $x \equiv_h y$, to $h(x) = h(y)$.

	Pokażemy, że $h_*$ jest homomorfizmem.
	Jeśli $f^n \in \Sigma$ jest $n$-arnym symbolem
		i odpowiadają mu w $A, B, A/_{\equiv_h}, \mathrm{Im}\,h$
		odpowiednio funkcje $f_A, f_B, f_{A/_{\equiv_h}}, f_{\mathrm{Im}\,h}$,
		to
	\begin{align}
		h_*(f_{A/_{\equiv_h}}([x_1]_{\equiv_h}, \dots, [x_n]_{\equiv_h})) &= \\
			(\text{definicja $A/_{\equiv_h}$ jako modelu})\qquad &= h_*([f_A(x_1, \dots, x_n)]_{\equiv_h}) \\
			(\text{definicja}\,h_*)\qquad &= h(f_A(x_1, \dots, x_n)) \\
			(\text{$h$ jest homomorfizmem})\qquad &= f_B(h(x_1), \dots, h(x_n)) \\
			(\,h(x_i) \in \mathrm{Im}\,h)\qquad &= f_{\mathrm{Im}\,h}(h(x_1), \dots, h(x_n)) \\
			(\text{definicja $h_*$})\qquad &= f_{\mathrm{Im}\,h}(h_*([x_1]_{\equiv_h}), \dots, h_*([x_n]_{\equiv_h})),
	\end{align}
	czyli $h_*$ jest homomorfizmem.

	Pozostaje pokazać,
		że $h_*$ jest bijekcją.
	Załóżmy, że $h_*([x]_{\equiv_h}) = h_*([y]_{\equiv_h})$.
	Czyli
	\begin{equation}
		h(x) = h_*([x]_{\equiv_h}) = h_*([y]_{\equiv_h}) = h(y).
	\end{equation}
	Więc $x \equiv_h y$ i tym samym $[x]_{\equiv_h} = [y]_{\equiv_h}$, czyli
		$h_*$ jest monomorfizmem.
	Załóżmy teraz,
		że $y \in \mathrm{Im}\,h$.
	Czyli istnieje $x \in A$, że $h(x) = y$.
	Mamy więc
	\begin{equation}
		h_*([x]_{\equiv_h}) = h(x) = y,
	\end{equation}
	czyli $h_*$ jest epimorfizmem.
	$h_*$ jest więc izomorfizmem.
\end{proof}

\subsection{Zadania.}

\begin{exercise}\label{ex:first-cong}
	Niech $\mathfrak{A} = (\Sigma, E)$ będzie dowolną algebrą.
	Niech $A$ będzie modelem $\mathfrak{A}$,
		a $\sim$ kongruencją na $A$.
	Udowodnij,
		że jeśli $F: A^n \to A$ jest funkcją zdefiniowaną przez wyrażenie składające się
		ze zmiennych i funkcji odpowiadającym wyrażeniom z $\Sigma$,
		to możemy zdefiniować $F_\sim : (A/_\sim)^n \to A/_\sim$ jako
	\begin{equation}
		F_\sim([x_1], \dots, [x_n]) = [F(x_1, \dots, x_n)],
	\end{equation}
		tzn., że definicja ta nie zależy od wyboru reprezentantów klas abstrakcji.

	\textit{Podpowiedź: indukcja po wysokości drzewa wyrażenia, które definiuje $F$.
	Przypadek bazowy jest już rozważony w dowodzie twierdzenia~\ref{th:cong-model}}
\end{exercise}

\begin{exercise}\label{ex:last-cong}
	Używając wyniku zadania~\ref{ex:first-cong} udowodnij brakującą część dowodu twierdzenia~\ref{th:cong-model}
\end{exercise}

\begin{exercise}\label{ex:ker-equiv}
	Udowodnij,
		że jeśli $h: A \to B$ jest homomorfizmem,
		to jego jądro $\equiv_h$ jest relacją równoważności
		(korzystając tylko z definicji jądra).
\end{exercise}

\begin{exercise}\label{ex:ker-submodel}
	Niech $h: A \to B$ będzie homomorfizmem dwóch modeli pewnej algebry.
	Udowodnij, że $\mathrm{Ker}\,h$ jest podmodelem $A$.
\end{exercise}

\begin{exercise}
	Sformułuj i udowodnij przypadek szczególny twierdzenia~\ref{th:isomorphism} dla monoidów.
\end{exercise}

% TODO: Dodać zadania ,,udowodnij, że ~ jest kongruencją"
% TODO: Dodać zadania ,,wyznacz jądro XXX"
